% GENERATED FILE. Edit canonical lessons, then run npm run generate:books.
% canonical-source-hash: fnv1a64:21a8d845f826b448
% canonical-lessons: RU-C416-provodnik, RU-C416-tambur, RU-C416-nosilshchik, RU-C416-telezhka, RU-C416-predyavlyat

\chapter{On board the train}
\label{ch:ru-on-board-the-train}
\hlchaptermodality{416}

\begin{chapteropening}
\textbf{By the end of this chapter:} \emph{I can manage on board a train: the attendant, the vestibule, the porter, the trolley and showing my ticket.}

Complete the last lesson of chapter 416: I can manage on board a train: the attendant, the vestibule, the porter, the trolley and showing my ticket.
\end{chapteropening}

\section[\texorpdfstring{\ru{проводник} (provodník)}{проводник (provodník)}]{\ru{проводник} (provodník) — a sleeping-car attendant; a guide}

\label{lesson:RU-C416-provodnik}



\begin{quote}
Before the new one: say the Russian for to announce, then the Russian for to make it.
\end{quote}

\subsection*{You'll want to know: \ru{проводник}}
\textbf{\ru{проводник}} (\emph{provodník}) — \textquotedblleft{}a sleeping-car attendant; a guide\textquotedblright{}.

\begin{quote}
\ru{Проводник} \ru{принёс} \ru{нам} \ru{чай}.
\end{quote}

\begin{quote}
\emph{Provodník prinyós nam chay.}
\end{quote}

\begin{quote}
The attendant brought us tea.
\end{quote}

Feminine \emph{\ru{проводница}}. The stress moves: \emph{\ru{проводника}} (\emph{provodniká}). Also a conductor of electricity.

\begin{cousinweb}
\emph{\ru{проводник}} is built on \emph{\ru{проводить}} 'to lead through' with \emph{-\ru{ник}} 'one who'. \emph{\ru{Водить}} is inherited from Proto-Slavic \emph{voditi} 'to lead' (EtymDB).
\end{cousinweb}

\begin{grammarlens}[title={What you've built}]
One more word for telling, explaining and describing at B1.
\end{grammarlens}

\subsection*{Your turn}
\begin{itemize}
\raggedright
  \item \emph{Say it:} \emph{provodník}
  \item \emph{Say it:} \emph{provodník}, once more
  \item \emph{Say it:} a sentence of your own with \emph{provodník}
  \item \emph{Recall:} say the Russian for to announce, then the Russian for to make it, then say \emph{provodník} again
\end{itemize}

\begin{tcolorbox}[breakable,colback=teal!4,colframe=teal!35!black,title={Before you move on}]
What does \ru{проводник} mean? (\textquotedblleft{}A sleeping-car attendant; a guide\textquotedblright{}.) Say it once more.
\end{tcolorbox}
\section[\texorpdfstring{\ru{тамбур} (támbur)}{тамбур (támbur)}]{\ru{тамбур} (támbur) — a vestibule (between train carriages)}

\label{lesson:RU-C416-tambur}



\begin{quote}
Before the new one: say the Russian for to make it, then the Russian for a sleeping-car attendant.
\end{quote}

\subsection*{You'll want to know: \ru{тамбур}}
\textbf{\ru{тамбур}} (\emph{támbur}) — \textquotedblleft{}a vestibule (between train carriages)\textquotedblright{}.

\begin{quote}
\ru{Пассажир} \ru{разговаривал} \ru{по} \ru{телефону} \ru{в} \ru{тамбуре}.
\end{quote}

\begin{quote}
\emph{Passazhír razgovárival po telefónu v támbure.}
\end{quote}

\begin{quote}
The passenger was talking on the phone in the vestibule.
\end{quote}

Masculine. Also the small space inside the door of a building.

\begin{cousinweb}
\emph{\ru{тамбур}} is usually traced to French \emph{tambour} 'a drum', which EtymDB traces to Arabic \emph{ṭunbūr}; the sense of a small closed space came later. EtymDB does not list the Russian word.
\end{cousinweb}

\begin{grammarlens}[title={What you've built}]
One more word for telling, explaining and describing at B1.
\end{grammarlens}

\subsection*{Your turn}
\begin{itemize}
\raggedright
  \item \emph{Say it:} \emph{támbur}
  \item \emph{Say it:} \emph{támbur}, once more
  \item \emph{Say it:} a sentence of your own with \emph{támbur}
  \item \emph{Recall:} say the Russian for to make it, then the Russian for a sleeping-car attendant, then say \emph{támbur} again
\end{itemize}

\begin{tcolorbox}[breakable,colback=teal!4,colframe=teal!35!black,title={Before you move on}]
What does \ru{тамбур} mean? (\textquotedblleft{}A vestibule (between train carriages)\textquotedblright{}.) Say it once more.
\end{tcolorbox}
\section[\texorpdfstring{\ru{носильщик} (nosílshchik)}{носильщик (nosílshchik)}]{\ru{носильщик} (nosílshchik) — a porter}

\label{lesson:RU-C416-nosilshchik}



\begin{quote}
Before the new one: say the Russian for a sleeping-car attendant, then the Russian for a vestibule.
\end{quote}

\subsection*{You'll want to know: \ru{носильщик}}
\textbf{\ru{носильщик}} (\emph{nosílshchik}) — \textquotedblleft{}a porter\textquotedblright{}.

\begin{quote}
\ru{Носильщик} \ru{помог} \ru{нам} \ru{с} \ru{чемоданами}.
\end{quote}

\begin{quote}
\emph{Nosílshchik pomóg nam s chemodánami.}
\end{quote}

\begin{quote}
The porter helped us with the suitcases.
\end{quote}

Masculine. Built on \emph{\ru{носить}} — to carry.

\begin{cousinweb}
\emph{\ru{носильщик}} is built on \emph{\ru{носить}} 'to carry' with \emph{-\ru{льщик}} 'one who does it'. \emph{\ru{Носить}} is inherited from Proto-Slavic \emph{nositi} (EtymDB).
\end{cousinweb}

\begin{grammarlens}[title={What you've built}]
One more word for telling, explaining and describing at B1.
\end{grammarlens}

\subsection*{Your turn}
\begin{itemize}
\raggedright
  \item \emph{Say it:} \emph{nosílshchik}
  \item \emph{Say it:} \emph{nosílshchik}, once more
  \item \emph{Say it:} a sentence of your own with \emph{nosílshchik}
  \item \emph{Recall:} say the Russian for a sleeping-car attendant, then the Russian for a vestibule, then say \emph{nosílshchik} again
\end{itemize}

\begin{tcolorbox}[breakable,colback=teal!4,colframe=teal!35!black,title={Before you move on}]
What does \ru{носильщик} mean? (\textquotedblleft{}A porter\textquotedblright{}.) Say it once more.
\end{tcolorbox}
\section[\texorpdfstring{\ru{тележка} (telézhka)}{тележка (telézhka)}]{\ru{тележка} (telézhka) — a trolley, a cart}

\label{lesson:RU-C416-telezhka}



\begin{quote}
Before the new one: say the Russian for a vestibule, then the Russian for a porter.
\end{quote}

\subsection*{You'll want to know: \ru{тележка}}
\textbf{\ru{тележка}} (\emph{telézhka}) — \textquotedblleft{}a trolley, a cart\textquotedblright{}.

\begin{quote}
\ru{На} \ru{вокзале} \ru{мы} \ru{взяли} \ru{тележку} \ru{для} \ru{багажа}.
\end{quote}

\begin{quote}
\emph{Na vokzále my vzyáli telézhku dlya bagázha.}
\end{quote}

\begin{quote}
At the station we took a trolley for the luggage.
\end{quote}

Feminine; a diminutive of \emph{\ru{телега}} — a cart. In shops too: \emph{\ru{тележка} \ru{для} \ru{покупок}}.

\begin{cousinweb}
\emph{\ru{тележка}} is a diminutive of \emph{\ru{телега}} 'a cart' (EtymDB lists \emph{\ru{тележка}}). The origin of \emph{\ru{телега}} itself is not settled.
\end{cousinweb}

\begin{grammarlens}[title={What you've built}]
One more word for telling, explaining and describing at B1.
\end{grammarlens}

\subsection*{Your turn}
\begin{itemize}
\raggedright
  \item \emph{Say it:} \emph{telézhka}
  \item \emph{Say it:} \emph{telézhka}, once more
  \item \emph{Say it:} a sentence of your own with \emph{telézhka}
  \item \emph{Recall:} say the Russian for a vestibule, then the Russian for a porter, then say \emph{telézhka} again
\end{itemize}

\begin{tcolorbox}[breakable,colback=teal!4,colframe=teal!35!black,title={Before you move on}]
What does \ru{тележка} mean? (\textquotedblleft{}A trolley, a cart\textquotedblright{}.) Say it once more.
\end{tcolorbox}
\section[\texorpdfstring{\ru{предъявлять} (predyavlyát)}{предъявлять (predyavlyát)}]{\ru{предъявлять} (predyavlyát) — to present (a ticket or a document)}

\label{lesson:RU-C416-predyavlyat}



\begin{quote}
Before the new one: say the Russian for a porter, then the Russian for a trolley.
\end{quote}

\subsection*{You'll want to know: \ru{предъявлять}}
\textbf{\ru{предъявлять}} (\emph{predyavlyát}) — \textquotedblleft{}to present (a ticket or a document)\textquotedblright{}.

\begin{quote}
\ru{В} \ru{вагоне} \ru{пассажиры} \ru{предъявляют} \ru{билеты} \ru{и} \ru{паспорта}.
\end{quote}

\begin{quote}
\emph{V vagóne passazhíry predyavlyáyut biléty i pasportá.}
\end{quote}

\begin{quote}
In the carriage the passengers present their tickets and passports.
\end{quote}

Imperfective; perfective \emph{\ru{предъявить}} (\emph{\ru{предъявлю}, \ru{предъявишь}}). Formal.

\begin{cousinweb}
\emph{\ru{предъявлять}} is built from \emph{\ru{пред}-} 'before' and \emph{\ru{являть}} 'to show': to put something before someone's eyes. EtymDB traces \emph{\ru{явить}} to Proto-Slavic \emph{aviti} 'to make visible'.
\end{cousinweb}

\begin{grammarlens}[title={What you've built}]
One more word for telling, explaining and describing at B1.
\end{grammarlens}

\subsection*{Your turn}
\begin{itemize}
\raggedright
  \item \emph{Say it:} \emph{predyavlyát}
  \item \emph{Say it:} \emph{predyavlyát}, once more
  \item \emph{Say it:} a sentence of your own with \emph{predyavlyát}
  \item \emph{Recall:} say the Russian for a porter, then the Russian for a trolley, then say \emph{predyavlyát} again
\end{itemize}

\begin{tcolorbox}[breakable,colback=teal!4,colframe=teal!35!black,title={Before you move on}]
What does \ru{предъявлять} mean? (\textquotedblleft{}To present (a ticket or a document)\textquotedblright{}.) Say it once more.
\end{tcolorbox}
